\documentclass[11pt]{article}
\usepackage[margin=1in]{geometry}
\usepackage{fancyhdr}
\usepackage{graphicx}
\usepackage{array}
\usepackage{booktabs}
\usepackage{xcolor}
\usepackage{tikz}
\usepackage{enumitem}

\pagestyle{fancy}
\fancyhf{}
\rhead{OFFICIAL USE ONLY}
\lhead{U.S. Department of Labor - OSHA}
\rfoot{Page \thepage}

\definecolor{oshablue}{RGB}{0,51,102}

\begin{document}

\begin{center}
{\Large\bfseries U.S. DEPARTMENT OF LABOR}\\
{\large Occupational Safety and Health Administration}\\[1em]
{\LARGE\bfseries INSPECTION REPORT}\\[0.5em]
{\large Inspection Number: 1654892}
\end{center}

\vspace{1em}
\hrule height 2pt
\vspace{1em}

\section*{ESTABLISHMENT INFORMATION}

\begin{tabular}{@{}p{4.5cm}p{10cm}@{}}
\textbf{Establishment:} & Titan Construction LLC \\
\textbf{DBA:} & Titan Builders \\
\textbf{Site Address:} & 450 Industrial Blvd, Riverside, CA 92507 \\
\textbf{Mailing Address:} & 1200 Commerce Way, Suite 300, Riverside, CA 92507 \\
\textbf{NAICS Code:} & 236220 - Commercial Building Construction \\
\textbf{Inspection Type:} & Accident Investigation \\
\textbf{Inspection Date:} & February 12, 2026 \\
\textbf{CSHO:} & Robert Martinez, Compliance Safety \& Health Officer \\
\end{tabular}

\section*{ACCIDENT SUMMARY}

On February 8, 2026, at approximately 10:15 AM, a 35-year-old female carpenter fell approximately 15 feet from scaffolding at a commercial construction site. The employee sustained a fractured wrist, concussion, and back injuries. The employee was hospitalized.

\section*{INSPECTION SCOPE}

This was an accident investigation inspection limited to the immediate work area and conditions related to the accident. The inspection included:

\begin{itemize}
\item Physical examination of the scaffolding involved in the accident
\item Review of safety training documentation
\item Interviews with employer representatives and witnesses
\item Review of company safety programs and policies
\item Photographic documentation of conditions
\end{itemize}

\section*{CITATIONS ISSUED}

\vspace{0.5em}
\noindent\fbox{\parbox{\textwidth}{
\textbf{Citation 1 - SERIOUS}\\[0.5em]
\textbf{29 CFR 1926.451(g)(1)(vii)} - Scaffolding - Guardrail Systems\\[0.5em]
The employer did not ensure that guardrail systems were installed along all open sides and ends of platforms. On or about February 8, 2026, at the Riverside Medical Plaza construction site, Building C, second floor, the east side of the scaffolding platform lacked guardrails, exposing employees to fall hazards of approximately 15 feet.\\[0.5em]
\textbf{Proposed Penalty: \$15,625}
}}

\vspace{1em}

\noindent\fbox{\parbox{\textwidth}{
\textbf{Citation 2 - SERIOUS}\\[0.5em]
\textbf{29 CFR 1926.502(d)} - Fall Protection - Duty to Have Fall Protection\\[0.5em]
The employer did not ensure that each employee on a walking/working surface with an unprotected side or edge which was 6 feet or more above a lower level was protected by the use of guardrail systems, safety net systems, or personal fall arrest systems. On February 8, 2026, employees working on scaffolding 15 feet above the floor were not provided with personal fall arrest systems (safety harnesses).\\[0.5em]
\textbf{Proposed Penalty: \$15,625}
}}

\vspace{1em}

\noindent\fbox{\parbox{\textwidth}{
\textbf{Citation 3 - SERIOUS}\\[0.5em]
\textbf{29 CFR 1926.21(b)(2)} - Safety Training and Education\\[0.5em]
The employer did not instruct each employee in the recognition and avoidance of unsafe conditions and the regulations applicable to the work environment to control or eliminate any hazards. Review of training records revealed that the injured employee had not received documented fall protection training within the past 12 months.\\[0.5em]
\textbf{Proposed Penalty: \$10,625}
}}

\vspace{1em}

\noindent\fbox{\parbox{\textwidth}{
\textbf{Citation 4 - SERIOUS}\\[0.5em]
\textbf{29 CFR 1926.451(f)(7)} - Scaffolding - Competent Person\\[0.5em]
A competent person was not present to inspect scaffolding for visible defects before each work shift. Records indicate the safety coordinator was off-site on the day of the accident, and no qualified substitute was assigned to perform required scaffolding inspections.\\[0.5em]
\textbf{Proposed Penalty: \$5,625}
}}

\vspace{1.5em}

\section*{PENALTY SUMMARY}

\begin{tabular}{@{}lr@{}}
\toprule
\textbf{Citation} & \textbf{Proposed Penalty} \\
\midrule
Citation 1 - Guardrail Systems & \$15,625 \\
Citation 2 - Fall Protection & \$15,625 \\
Citation 3 - Safety Training & \$10,625 \\
Citation 4 - Competent Person & \$5,625 \\
\midrule
\textbf{TOTAL PROPOSED PENALTIES} & \textbf{\$47,500} \\
\bottomrule
\end{tabular}

\section*{ABATEMENT REQUIREMENTS}

The employer must abate the cited hazards by the following dates:

\begin{enumerate}
\item Citation 1: Immediate upon receipt
\item Citation 2: Immediate upon receipt
\item Citation 3: Within 15 calendar days
\item Citation 4: Within 7 calendar days
\end{enumerate}

\section*{EMPLOYER RIGHTS}

The employer has 15 working days from receipt of these citations to:
\begin{itemize}
\item Pay the penalty
\item Request an informal conference with the OSHA Area Director
\item Contest the citations and/or penalties before the Occupational Safety and Health Review Commission
\end{itemize}

\vspace{2em}

\begin{tabular}{@{}p{7cm}p{7cm}@{}}
\hrulefill & \\
\textbf{Robert Martinez} & \textbf{Date:} February 22, 2026 \\
Compliance Safety \& Health Officer & \\
OSHA - Region IX & \\
\end{tabular}

\vspace{1em}
\begin{center}
\small\textit{Citation Date: February 22, 2026}
\end{center}

\end{document}
